% GENERATED FILE. Edit canonical lessons, then run npm run generate:books.
% canonical-source-hash: fnv1a64:b9bc4fbaeab299c9
% canonical-lessons: RU-C274-postepenno, RU-C274-odnovremenno, RU-C274-ponachalu, RU-C274-vposledstvii, RU-C274-naposledok

\chapter{Little by little}
\label{ch:ru-little-by-little}
\hlchaptermodality{274}

\begin{chapteropening}
\textbf{By the end of this chapter:} \emph{I can say how events unfolded: gradually, at the same time, at first, later on and at the very end.}

Complete the last lesson of chapter 274: I can say how events unfolded: gradually, at the same time, at first, later on and at the very end.
\end{chapteropening}

\section[\texorpdfstring{\ru{постепенно} (postépenno)}{постепенно (postépenno)}]{\ru{постепенно} (postépenno) — gradually, little by little}

\label{lesson:RU-C274-postepenno}



\begin{quote}
Before the new one: say the Russian for to stop, then the Russian for to come.
\end{quote}

\subsection*{You'll want to know: \ru{постепенно}}
\textbf{\ru{постепенно}} (\emph{postépenno}) — \textquotedblleft{}gradually, little by little\textquotedblright{}.

\begin{quote}
\ru{Постепенно} \ru{все} \ru{поняли} \ru{правила} \ru{игры}.
\end{quote}

\begin{quote}
\emph{Postépenno vse pónyali právila igrý.}
\end{quote}

\begin{quote}
Gradually everyone understood the rules of the game.
\end{quote}

The adverb of \emph{\ru{постепенный}} (\emph{postépennyy}) \textquotedblleft{}gradual\textquotedblright{}.

\begin{cousinweb}
\emph{\ru{постепенно}} is built on \emph{\ru{по}} and \emph{\ru{степень}} \textquotedblleft{}a degree, a step\textquotedblright{}, inherited from Proto-Slavic \emph{stepenĭ} (EtymDB): step by step.
\end{cousinweb}

\begin{grammarlens}[title={What you've built}]
One more word for telling, explaining and describing at B1.
\end{grammarlens}

\subsection*{Your turn}
\begin{itemize}
\raggedright
  \item \emph{Say it:} \emph{postépenno}
  \item \emph{Say it:} \emph{postépenno}, once more
  \item \emph{Say it:} a sentence of your own with \emph{postépenno}
  \item \emph{Recall:} say the Russian for to stop, then the Russian for to come, then say \emph{postépenno} again
\end{itemize}

\begin{tcolorbox}[breakable,colback=teal!4,colframe=teal!35!black,title={Before you move on}]
What does \ru{постепенно} mean? (\textquotedblleft{}Gradually, little by little\textquotedblright{}.) Say it once more.
\end{tcolorbox}
\section[\texorpdfstring{\ru{одновременно} (odnovreménno)}{одновременно (odnovreménno)}]{\ru{одновременно} (odnovreménno) — at the same time, simultaneously}

\label{lesson:RU-C274-odnovremenno}



\begin{quote}
Before the new one: say the Russian for to come, then the Russian for gradually.
\end{quote}

\subsection*{You'll want to know: \ru{одновременно}}
\textbf{\ru{одновременно}} (\emph{odnovreménno}) — \textquotedblleft{}at the same time, simultaneously\textquotedblright{}.

\begin{quote}
\ru{Они} \ru{пришли} \ru{одновременно}.
\end{quote}

\begin{quote}
\emph{Oní prishlí odnovreménno.}
\end{quote}

\begin{quote}
They arrived at the same time.
\end{quote}

\emph{\ru{Одновременно} \ru{с}} + instrumental means \textquotedblleft{}at the same time as\textquotedblright{}: \emph{\ru{одновременно} \ru{с} \ru{нами}}.

\begin{cousinweb}
\emph{\ru{одновременно}} joins \emph{\ru{один}} \textquotedblleft{}one\textquotedblright{}, inherited from Proto-Slavic \emph{edinŭ} (EtymDB), and \emph{\ru{время}} \textquotedblleft{}time\textquotedblright{}, usually taken as a Church Slavonic shape beside Old Russian \emph{\ru{веремя}}; EtymDB lists \emph{\ru{время}} with no parent.
\end{cousinweb}

\begin{grammarlens}[title={What you've built}]
One more word for telling, explaining and describing at B1.
\end{grammarlens}

\subsection*{Your turn}
\begin{itemize}
\raggedright
  \item \emph{Say it:} \emph{odnovreménno}
  \item \emph{Say it:} \emph{odnovreménno}, once more
  \item \emph{Say it:} a sentence of your own with \emph{odnovreménno}
  \item \emph{Recall:} say the Russian for to come, then the Russian for gradually, then say \emph{odnovreménno} again
\end{itemize}

\begin{tcolorbox}[breakable,colback=teal!4,colframe=teal!35!black,title={Before you move on}]
What does \ru{одновременно} mean? (\textquotedblleft{}At the same time, simultaneously\textquotedblright{}.) Say it once more.
\end{tcolorbox}
\section[\texorpdfstring{\ru{поначалу} (ponachálu)}{поначалу (ponachálu)}]{\ru{поначалу} (ponachálu) — at first}

\label{lesson:RU-C274-ponachalu}



\begin{quote}
Before the new one: say the Russian for gradually, then the Russian for at the same time.
\end{quote}

\subsection*{You'll want to know: \ru{поначалу}}
\textbf{\ru{поначалу}} (\emph{ponachálu}) — \textquotedblleft{}at first\textquotedblright{}.

\begin{quote}
\ru{Поначалу} \ru{мне} \ru{было} \ru{трудно}.
\end{quote}

\begin{quote}
\emph{Ponachálu mne býlo trúdno.}
\end{quote}

\begin{quote}
At first it was hard for me.
\end{quote}

Colloquial, and it hints that things changed later. The neutral word is \emph{\ru{сначала}} (\emph{snachála}).

\begin{cousinweb}
\emph{\ru{поначалу}} is built on \emph{\ru{по}} and \emph{\ru{начало}} \textquotedblleft{}beginning\textquotedblright{}, inherited through Old East Slavic \emph{načęlo} from Proto-Slavic \emph{načędlo} (EtymDB), the root of \emph{\ru{начать}} \textquotedblleft{}to begin\textquotedblright{}.
\end{cousinweb}

\begin{grammarlens}[title={What you've built}]
One more word for telling, explaining and describing at B1.
\end{grammarlens}

\subsection*{Your turn}
\begin{itemize}
\raggedright
  \item \emph{Say it:} \emph{ponachálu}
  \item \emph{Say it:} \emph{ponachálu}, once more
  \item \emph{Say it:} a sentence of your own with \emph{ponachálu}
  \item \emph{Recall:} say the Russian for gradually, then the Russian for at the same time, then say \emph{ponachálu} again
\end{itemize}

\begin{tcolorbox}[breakable,colback=teal!4,colframe=teal!35!black,title={Before you move on}]
What does \ru{поначалу} mean? (\textquotedblleft{}At first\textquotedblright{}.) Say it once more.
\end{tcolorbox}
\section[\texorpdfstring{\ru{впоследствии} (vposlédstvii)}{впоследствии (vposlédstvii)}]{\ru{впоследствии} (vposlédstvii) — later on, subsequently}

\label{lesson:RU-C274-vposledstvii}



\begin{quote}
Before the new one: say the Russian for at the same time, then the Russian for at first.
\end{quote}

\subsection*{You'll want to know: \ru{впоследствии}}
\textbf{\ru{впоследствии}} (\emph{vposlédstvii}) — \textquotedblleft{}later on, subsequently\textquotedblright{}.

\begin{quote}
\ru{Впоследствии} \ru{он} \ru{стал} \ru{известным} \ru{врачом}.
\end{quote}

\begin{quote}
\emph{Vposlédstvii on stal izvéstnym vrachóm.}
\end{quote}

\begin{quote}
Later on he became a well-known doctor.
\end{quote}

Formal and written. In speech: \emph{\ru{потом}}, \emph{\ru{позже}}.

\begin{cousinweb}
\emph{\ru{впоследствии}} is built on \emph{\ru{в}} and \emph{\ru{последствие}} \textquotedblleft{}a consequence\textquotedblright{}, from \emph{\ru{след}} \textquotedblleft{}a track, a trace\textquotedblright{} (EtymDB): what comes in the footsteps of an event.
\end{cousinweb}

\begin{grammarlens}[title={What you've built}]
One more word for telling, explaining and describing at B1.
\end{grammarlens}

\subsection*{Your turn}
\begin{itemize}
\raggedright
  \item \emph{Say it:} \emph{vposlédstvii}
  \item \emph{Say it:} \emph{vposlédstvii}, once more
  \item \emph{Say it:} a sentence of your own with \emph{vposlédstvii}
  \item \emph{Recall:} say the Russian for at the same time, then the Russian for at first, then say \emph{vposlédstvii} again
\end{itemize}

\begin{tcolorbox}[breakable,colback=teal!4,colframe=teal!35!black,title={Before you move on}]
What does \ru{впоследствии} mean? (\textquotedblleft{}Later on, subsequently\textquotedblright{}.) Say it once more.
\end{tcolorbox}
\section[\texorpdfstring{\ru{напоследок} (naposlédok)}{напоследок (naposlédok)}]{\ru{напоследок} (naposlédok) — at the very end, last of all}

\label{lesson:RU-C274-naposledok}



\begin{quote}
Before the new one: say the Russian for at first, then the Russian for later on.
\end{quote}

\subsection*{You'll want to know: \ru{напоследок}}
\textbf{\ru{напоследок}} (\emph{naposlédok}) — \textquotedblleft{}at the very end, last of all\textquotedblright{}.

\begin{quote}
\ru{Напоследок} \ru{он} \ru{улыбнулся} \ru{и} \ru{ушёл}.
\end{quote}

\begin{quote}
\emph{Naposlédok on ulybnúlsya i ushyól.}
\end{quote}

\begin{quote}
At the very end he smiled and left.
\end{quote}

Colloquial: the last thing before parting.

\begin{cousinweb}
\emph{\ru{напоследок}} is built on \emph{\ru{на}} and \emph{\ru{последок}} \textquotedblleft{}the last bit\textquotedblright{}, from \emph{\ru{последний}} \textquotedblleft{}last\textquotedblright{}, which goes back to \emph{\ru{след}} \textquotedblleft{}a track\textquotedblright{} (EtymDB).
\end{cousinweb}

\begin{grammarlens}[title={What you've built}]
One more word for telling, explaining and describing at B1.
\end{grammarlens}

\subsection*{Your turn}
\begin{itemize}
\raggedright
  \item \emph{Say it:} \emph{naposlédok}
  \item \emph{Say it:} \emph{naposlédok}, once more
  \item \emph{Say it:} a sentence of your own with \emph{naposlédok}
  \item \emph{Recall:} say the Russian for at first, then the Russian for later on, then say \emph{naposlédok} again
\end{itemize}

\begin{tcolorbox}[breakable,colback=teal!4,colframe=teal!35!black,title={Before you move on}]
What does \ru{напоследок} mean? (\textquotedblleft{}At the very end, last of all\textquotedblright{}.) Say it once more.
\end{tcolorbox}
